% ============================================================================
% ARBEITSBLATT DS 6 - Gruppe B: Eigenschaften für Berufe (Brückenthema)
% ============================================================================

\documentclass[11pt,a4paper]{article}

\usepackage[utf8]{inputenc}
\usepackage[T1]{fontenc}
\usepackage[ngerman]{babel}
\usepackage[left=2cm,right=2cm,top=1.8cm,bottom=1.5cm]{geometry}
\usepackage{helvet}
\usepackage[table]{xcolor}
\usepackage{tabularx}
\usepackage{array}
\usepackage{enumitem}
\usepackage{tcolorbox}
\usepackage{fancyhdr}
\usepackage{lastpage}
\usepackage{amssymb}

\renewcommand{\familydefault}{\sfdefault}

\definecolor{spanisch}{HTML}{c41e3a}
\definecolor{grau}{HTML}{666666}
\definecolor{hellgrau}{HTML}{f5f5f5}
\definecolor{gruppeB}{HTML}{1565c0}

\pagestyle{fancy}
\fancyhf{}
\fancyhead[L]{\small\textcolor{grau}{Spanisch Spätbeginner -- Gruppe B}}
\fancyhead[R]{\small\textcolor{grau}{AB-06-B}}
\fancyfoot[C]{\small\thepage/\pageref{LastPage}}
\renewcommand{\headrulewidth}{0.5pt}

\newcommand{\luecke}[1]{\underline{\hspace{#1}}}
\newcommand{\punktebox}[1]{\hfill\fbox{\small \_\_/#1}}
\newcommand{\es}[1]{\textit{#1}}

\newtcolorbox{merkkasten}[1][]{
    colback=white,
    colframe=gruppeB,
    fonttitle=\bfseries\color{gruppeB},
    title=#1,
    boxrule=1.5pt,
    arc=0pt,
    left=5pt, right=5pt, top=5pt, bottom=5pt
}

\newtcolorbox{buchverweis}{
    colback=hellgrau,
    colframe=grau,
    boxrule=0.5pt,
    arc=0pt,
    left=5pt, right=5pt, top=3pt, bottom=3pt
}

\newcounter{aufgabe}
\newcommand{\aufgabe}[2]{%
    \stepcounter{aufgabe}%
    \vspace{0.8em}%
    \noindent\textbf{\theaufgabe.} #1 \punktebox{#2}%
    \vspace{0.3em}%
}

\begin{document}

% ============================================================================
% KOPFBEREICH
% ============================================================================
\noindent
\begin{tabularx}{\textwidth}{@{}Xr@{}}
    {\Large\textbf{\textcolor{gruppeB}{Profesiones y características}}} & Name: \luecke{5cm} \\[0.3em]
    {\small DS 6 -- Brückenthema} & Datum: \luecke{3cm} \\
\end{tabularx}

\vspace{0.3em}
\hrule
\vspace{0.8em}

% ============================================================================
% MERKKASTEN 1: hay que + Infinitiv
% ============================================================================
\begin{merkkasten}[Kasten 1: hay que + Infinitiv (Notwendigkeit)]
\vspace{0.2em}

\begin{tabularx}{\textwidth}{@{}|p{4cm}|X|@{}}
\hline
\rowcolor{hellgrau}
\textbf{Struktur} & \textbf{Beispiel} \\
\hline
\textbf{hay que ser} + Adj. & Para ser médico, \textbf{hay que ser} paciente. \\[0.4em]
\hline
\textbf{hay que tener} + Subst. & Para ser profesor, \textbf{hay que tener} paciencia. \\[0.4em]
\hline
\textbf{hay que + Infinitiv} & Hay que \luecke{2.5cm} mucho. (arbeiten) \\[0.4em]
\hline
\end{tabularx}

\vspace{0.3em}
\textbf{Merkhilfe:} \es{hay que} = man muss (unpersönlich)
\vspace{0.2em}
\end{merkkasten}

\vspace{0.8em}

% ============================================================================
% MERKKASTEN 2: Adjektive für Eigenschaften
% ============================================================================
\begin{merkkasten}[Kasten 2: Adjektive zur Personenbeschreibung]
\vspace{0.2em}

\begin{tabularx}{\textwidth}{@{}|X|X|X|@{}}
\hline
\rowcolor{hellgrau}
\textbf{Positiv} & \textbf{Neutral} & \textbf{Kritisch} \\
\hline
simpático/-a (nett) & tranquilo/-a (ruhig) & perezoso/-a (faul) \\[0.3em]
\hline
paciente (geduldig) & serio/-a (ernst) & impaciente (ungeduldig) \\[0.3em]
\hline
responsable (verantwortungsvoll) & reservado/-a (zurückhaltend) & egoísta (egoistisch) \\[0.3em]
\hline
trabajador/-a (fleißig) & organizado/-a (organisiert) & -- \\[0.3em]
\hline
\end{tabularx}

\vspace{0.3em}
\end{merkkasten}

\vspace{1em}

% ============================================================================
% AUFGABEN
% ============================================================================

\aufgabe{Wortschatz: Berufe und passende Eigenschaften. Übersetze und ordne zu.}{8}

\begin{tabularx}{\textwidth}{@{}|X|X|X|@{}}
\hline
\rowcolor{hellgrau}
\textbf{Beruf (Spanisch)} & \textbf{Deutsch} & \textbf{Passende Eigenschaft} \\
\hline
el/la médico/-a & \luecke{2.5cm} & \luecke{3cm} \\[0.5em]
\hline
el/la profesor/-a & \luecke{2.5cm} & \luecke{3cm} \\[0.5em]
\hline
el/la enfermero/-a & \luecke{2.5cm} & \luecke{3cm} \\[0.5em]
\hline
el/la jefe/-a & \luecke{2.5cm} & \luecke{3cm} \\[0.5em]
\hline
\end{tabularx}

\vspace{0.8em}

\aufgabe{Bilde Sätze mit \es{hay que ser ... porque ...}}{8}

\begin{buchverweis}
\textbf{Beispiel:} \es{Para ser médico hay que ser paciente porque trabajas con personas enfermas.}
\end{buchverweis}

\begin{enumerate}[label=\alph*), leftmargin=2em]
    \item Para ser profesor hay que ser \luecke{2.5cm} porque \luecke{5cm}.
    \vspace{0.3cm}
    \item Para ser enfermero hay que ser \luecke{2.5cm} porque \luecke{5cm}.
    \vspace{0.3cm}
    \item Para ser jefe hay que ser \luecke{2.5cm} porque \luecke{5cm}.
    \vspace{0.3cm}
    \item Para ser \luecke{2.5cm} hay que ser \luecke{2.5cm} porque \luecke{4cm}.
\end{enumerate}

\vspace{0.8em}

\aufgabe{Mini-Diskussion: ¿Cómo es un buen jefe? Schreibe 3 Sätze.}{6}

\begin{buchverweis}
\textbf{Nutze:} \es{Un buen jefe es ... / tiene que ser ... / además es ... / pero no es ...}
\end{buchverweis}

\vspace{0.5em}
\begin{tabularx}{\textwidth}{@{}|X|@{}}
\hline
\\[0.5em]
1. \luecke{13cm} \\[0.8em]
2. \luecke{13cm} \\[0.8em]
3. \luecke{13cm} \\[0.5em]
\hline
\end{tabularx}

\vspace{1em}

\aufgabe{Overlap-Vorbereitung: Interview-Fragen üben.}{6}

\begin{buchverweis}
\textbf{Du wirst einen Schüler aus Gruppe A interviewen.} Übe diese Fragen:
\end{buchverweis}

\vspace{0.5em}
\textbf{Übersetze die Fragen ins Spanische:}

\vspace{0.3em}
a) Wie bist du? $\rightarrow$ \luecke{8cm}

\vspace{0.3cm}
b) Wie ist dein bester Freund / deine beste Freundin? $\rightarrow$ \luecke{6cm}

\vspace{0.3cm}
c) Wie ist deine Familie? $\rightarrow$ \luecke{8cm}

\vspace{0.3cm}
d) Warum? $\rightarrow$ \luecke{8cm}

\vspace{1em}

\aufgabe{Notizvorlage für das Interview (Overlap).}{5}

\begin{buchverweis}
\textbf{Notiere die Antworten deines Partners in Stichpunkten.}
\end{buchverweis}

\vspace{0.5em}
\textbf{Interview mit:} \luecke{4cm}

\vspace{0.3cm}
\begin{tabularx}{\textwidth}{@{}|X|@{}}
\hline
\rowcolor{hellgrau}
\textbf{Er/Sie ist:} \\
\hline
\\[0.8em]
\luecke{13cm} \\[0.5em]
\hline
\rowcolor{hellgrau}
\textbf{Sein/Ihr bester Freund ist:} \\
\hline
\\[0.8em]
\luecke{13cm} \\[0.5em]
\hline
\rowcolor{hellgrau}
\textbf{Seine/Ihre Familie ist:} \\
\hline
\\[0.8em]
\luecke{13cm} \\[0.5em]
\hline
\end{tabularx}

\vspace{0.5em}
\textbf{Zusammenfassung (auf Spanisch):}

\vspace{0.3cm}
\luecke{3cm} es \luecke{3cm} y \luecke{3cm}. También es \luecke{3cm}.

\vspace{1em}
\hrule
\vspace{0.3em}
\begin{flushright}
\textbf{Gesamt: \luecke{1cm} / 33 Punkte}
\end{flushright}

\end{document}
